%======================================================================
\chapter{From continuum to discrete: numerical methods}\label{ch:nc}
%======================================================================
% Labels of this chapter: *:nc-*  (bibitem keys nc-*). Status: outline, see ch03_continuation/README.md.

\framepart{Outline}
Chapters~\ref{ch:st} and~\ref{ch:ld} worked with a few unknowns, and every
derivative was taken by hand. A real structure is a continuum. Its potential
energy $\mathcal{E}(\vu,\eta)$ is an integral over the body, and the unknown is a
displacement field. This chapter turns the field into a column of numbers, and
shows that nothing of Chapter~\ref{ch:ld} is lost on the way. The first variation of
$\mathcal{E}$ is the weak form of equilibrium, the principle of virtual work. The
second variation is the stability operator. Finite elements replace the field by
nodal values $\mat{u}$, and the two variations become the residual vector and the
tangent stiffness matrix $\mat{K}$. The energy criterion, the critical point and the
mode are then exactly those of Chapter~\ref{ch:ld}, with many more unknowns.

Two kinds of computation follow. \emph{Linear buckling} assumes a linear prebuckling
state and asks for the load at which $\mat{K}$ becomes singular: a generalised
eigenvalue problem. Its smallest eigenvalue is the minimum of a Rayleigh quotient,
so the eigenproblem is itself a weak (variational) problem, and any finite element
or Ritz approximation gives an upper bound. \emph{Nonlinear continuation} follows the
equilibrium path with the arc-length methods of Riks and Crisfield, which are
Box~1.2 written for a residual vector and a load vector. Along the
path we detect the critical points, tell limit points from bifurcation points, and
switch to the new branch. The von Mises truss of Chapter~\ref{ch:ld} is the test case
throughout: two bar elements, closed-form answers, and two lessons on when linear
buckling is wrong.

\paragraph{By the end of this chapter you should be able to:}
\begin{itemize}
  \item write the weak form of equilibrium and the second variation of the energy of an
        elastic body, and recognise the initial-stress (geometric) term;
  \item derive the residual and the tangent stiffness of a finite element model from
        its energy, and split $\mat{K}$ into its material and geometric parts;
  \item set up linear buckling as the eigenproblem $(\mat{K}_0 + \eta\,\mat{K}_\sigma)\,{}^1\mat{u} = \mat{0}$,
        and say when it is (and is not) a good estimate;
  \item state the Rayleigh quotient and the min--max principle, and explain why a
        Ritz or finite element buckling load is an upper bound;
  \item follow a path through a limit point with the Riks--Crisfield arc-length method;
  \item detect a critical point from the eigenvalues or the inertia of $\mat{K}$,
        classify it, and switch to the bifurcated branch;
  \item run a first linear buckling and a first arc-length computation in FEniCS and in
        Cast3M, and check them against a closed form.
\end{itemize}

\paragraph{Plan of the chapter.}
Section~\ref{sec:nc-weak} goes from the energy of an elastic body to the weak form and
the second variation. Section~\ref{sec:nc-fe} discretises them by finite elements.
Section~\ref{sec:nc-lba} sets up linear buckling as an eigenproblem, and
Section~\ref{sec:nc-rayleigh} reads that eigenproblem as a minimisation.
Section~\ref{sec:nc-riks} follows nonlinear paths by arc-length continuation.
Section~\ref{sec:nc-detect} detects and classifies critical points, and
Section~\ref{sec:nc-switch} switches branches. Section~\ref{sec:nc-tools} presents the
two finite element codes. A summary, exercises (page~\pageref{sec:nc-exo}) and the
references close the chapter.

\draftnote{Open (thread, 21 Nov 2026): the planned section on the asymptotic numerical
method (Cochelin, Damil, Potier-Ferry) is not in the author's list. Proposal: one
paragraph at the end of Section~\ref{sec:nc-riks} as an alternative to
predictor--corrector schemes, or drop it. To be decided.}

%======================================================================
\section{From the energy to the weak form}\label{sec:nc-weak}
\index{weak form}\index{virtual work}%
\draftnote{Planned content:
(1)~Start small: the axially compressed column, energy
$\tfrac12\int EI\,w''^2\dd x - \tfrac12 P\int w'^2\dd x$; its first variation is the weak form,
integration by parts gives $EIw'''' + Pw'' = 0$ and the boundary conditions. Closed form
$P_c = \pi^2EI/L^2$ (the elastica proper is Chapter~\ref{ch:rd}).
(2)~Elastic body (notation of Constantinescu--Korsunsky): $\mathcal{E}(\vu,\eta) = \int_\Omega W(\tens{E})\dd V
- \eta\,\ell(\vu)$ over $\Kad$, Green--Lagrange strain $\tens{E}$, $W = \tfrac12\,\tens{E}:\bbC:\tens{E}$.
(3)~First variation $\mathcal{E}_{,\vu}[\delta\vu] = 0$ for all $\delta\vu \in \Kad_0$: the principle of
virtual work, the weak form of equilibrium.
(4)~Second variation $\mathcal{E}_{,\vu\vu}[\delta\vu,\delta\vu]$ = material part + initial-stress part
$\int \tens{S}:(\operatorname{grad}\delta\vu\T\operatorname{grad}\delta\vu)\dd V$; stability = positive second
variation (Trefftz criterion), the continuum version of the gbox of Section~\ref{sec:ld-energycrit}.
gbox: the weak form is the gradient of the energy, the stability operator is its Hessian.}

%======================================================================
\section{From fields to matrices: finite elements}\label{sec:nc-fe}
\index{finite elements}\index{tangent stiffness}%
\draftnote{Planned content:
(1)~Interpolation $\vu \approx \mat{N}\mat{u}$; $\mathcal{E}(\mat{u},\eta)$ becomes a function of the nodal
values: the discrete systems of Chapter~\ref{ch:ld} are recovered.
(2)~Residual $\mat{r}(\mat{u},\eta) = \mathcal{E}_{,\mat{u}} = \mat{f}_{\mathrm{int}}(\mat{u}) - \eta\,\mat{p}$
(reference load vector $\mat{p}$); tangent stiffness $\mat{K} = \mathcal{E}_{,\mat{u}\mat{u}}$
= material part (with initial displacements) + geometric part $\mat{K}_\sigma$.
(3)~Worked example: the two-bar truss of Section~\ref{sec:ld-snap} as two bar elements with
Green strain; assembly by hand gives exactly~\eqref{eq:ld-truss-K}.
(4)~Automatic differentiation: in FEniCS the residual and the tangent are derived from the
energy form (\texttt{derivative}); in Cast3M they are assembled element by element.
Box~3.1: from energy to residual and tangent.}

%======================================================================
\section{Linear buckling as an eigenproblem}\label{sec:nc-lba}
\index{linear buckling}\index{eigenvalue problem!buckling}%
\draftnote{Planned content:
(1)~Linear prebuckling: $\mat{K}_0\,{}^0\mat{u} = \eta\,\mat{p}$, stresses proportional to $\eta$, so
$\mat{K} \approx \mat{K}_0 + \eta\,\mat{K}_\sigma$.
(2)~Critical load: $(\mat{K}_0 + \eta\,\mat{K}_\sigma)\,{}^1\mat{u} = \mat{0}$, a generalised
symmetric eigenproblem; the smallest positive $\eta$ and its mode. Box~3.2.
(3)~Two warnings on the truss (checked with SymPy, $EA = L = 1$):
$\varphi = 30^\circ$, linear buckling gives $0.250$ for the true limit load $0.0481$ (a factor $5.2$:
the prebuckling is not linear, a snap-through is not a bifurcation);
$\varphi = 65^\circ$, it gives $0.324$ for the true sideways bifurcation at $0.2434$ ($+33\%$).
gbox: linear buckling is exact only when the prebuckling state is linear and the critical point
is a bifurcation (Bushnell 1981, Section~4).
(4)~Consistent linearisation $\mat{K}_0 + \eta\,d\mat{K}/d\eta$ (initial-displacement matrix included):
$0.0833$ at $30^\circ$, still too high.}

%======================================================================
\section{The eigenproblem as a minimisation}\label{sec:nc-rayleigh}
\index{Rayleigh quotient}\index{Rayleigh--Ritz method}\index{min--max principle}%
\draftnote{Planned content:
(1)~Rayleigh quotient: $\eta_c = \min\, a(\vv,\vv)/b(\vv,\vv)$
over $\Kad_0$ with $b > 0$, $a$ the material and $-b$ the initial-stress bilinear forms; for the
column $P_c = \min \int EI\,w''^2/\int w'^2$. The weak form of the eigenproblem is the
stationarity of this quotient.
(2)~Courant--Fischer min--max for the higher modes; orthogonality of modes in $a$.
(3)~Rayleigh--Ritz: a smaller space gives a larger minimum, so every Ritz or conforming finite
element buckling load is an upper bound. Column with $w = x(L - x)$: $12\,EI/L^2$ against
$\pi^2EI/L^2$ ($+21.6\%$); Hermite cubic elements converge from above.
(4)~Same structure as the free vibrations, $\omega^2 = $ Rayleigh quotient with the mass
(Section~\ref{sec:ld-energycrit}; Constantinescu--Korsunsky Section~7.5).
gbox: a discrete buckling load is never below the exact one (for conforming elements and a
fixed prebuckling stress).}

%======================================================================
\section{Following nonlinear paths: arc-length methods}\label{sec:nc-riks}
\index{continuation!arc-length}\index{Riks method}\index{Crisfield method}%
\draftnote{Planned content:
(1)~Load control fails at a limit point, displacement control fails at a snap-back;
Box~1.2 with $\mat{f} = \mat{r}$, $\mat{J} = \mat{K}$, $\mat{f}_r = -\mat{p}$.
(2)~Riks (linearised constraint, normal plane) and Crisfield (spherical constraint, quadratic
equation for $\delta\eta$, choice of the root); the bordering algorithm: two solves with
$\mat{K}$ per iteration, no augmented matrix needed (essential with a sparse FE solver).
(3)~Scaling of load and displacement in the arc length; step control by the number of
Newton iterations.
(4)~Test: the truss at $30^\circ$ through both limit points, against~\eqref{eq:ld-truss-principal}.
Box~3.3: Riks--Crisfield arc-length step.}

%======================================================================
\section{Detecting and classifying critical points}\label{sec:nc-detect}
\index{critical point!detection}\index{inertia of a matrix}%
\draftnote{Planned content:
(1)~Monitors along the path: sign of $\det\mat{K}$, smallest eigenvalues of $\mat{K}$
(SLEPc, shift-invert), and the inertia: the number of negative pivots of
$\mat{K} = \mat{L}\mat{D}\mat{L}\T$ counts the negative eigenvalues (Sylvester), at no extra cost
in a direct solver.
(2)~Locating the point: bisection on the step, or the extended system
$\mat{r} = \mat{0}$, $\mat{K}{}^1\mat{u} = \mat{0}$, $\norm{{}^1\mat{u}} = 1$ solved by Newton.
(3)~Classification by the test~\eqref{eq:ld-limit-test}: ${}^1\mat{u}\T\mat{p} \neq 0$ limit point,
$= 0$ bifurcation; the sign of the load component of the tangent.
(4)~Truss at $65^\circ$: one negative pivot appears at $0.2434$ with $\eta$ still increasing:
a bifurcation, mode $(1,0)$.}

%======================================================================
\section{Branch switching}\label{sec:nc-switch}
\index{branch switching}%
\draftnote{Planned content:
(1)~Directions of the branches from the second-order (bifurcation) equation, as in
Section~\ref{sec:st-continuation}; Koiter's $\eta_1$, $\eta_2$ (Box~2.2) give
the start of the new branch.
(2)~Practice in FE codes: predictor $\mat{u}_c + \epsilon\,{}^1\mat{u}$ followed by the corrector,
or an imperfection in the shape of the mode (load or geometry) and continuation of the
imperfect structure.
(3)~Truss at $65^\circ$: the bifurcated branch is the circle~\eqref{eq:ld-truss-branch}; with a
sideways load the maximum loads reproduce the two-thirds law of Section~\ref{sec:ld-powerlaws}.
gbox: an imperfection in the shape of the first mode is the safe way to switch, and it also
measures the imperfection sensitivity.}

%======================================================================
\section{Two finite element codes}\label{sec:nc-tools}
\index{FEniCS}\index{Cast3M}%
\draftnote{Planned content (short, a map rather than a manual):
FEniCS: the energy written in UFL, residual and tangent by \texttt{derivative}, Newton by PETSc
SNES, the eigenproblem by SLEPc (generalised Hermitian, shift-and-invert), arc-length written on
top of the Newton solver by bordering.
Cast3M: bar and beam elements, \texttt{KSIG} for the initial-stress matrix and \texttt{VIBR} for the
eigenproblem, \texttt{PASAPAS} for the nonlinear path with its arc-length option.
The same checks for both: closed form first, then mesh refinement.}

%======================================================================
\framepart{Summary}
\begin{gbox*}
\begin{itemize}
  \item \textbf{To be written.} \draftnote{take-home messages, one per section.}
\end{itemize}
\end{gbox*}

%======================================================================
\framepart{Exercises}\label{sec:nc-exo}
\framesub{Short questions}
\draftnote{Planned: Rayleigh quotient of the column with a trial function; which is larger,
the Ritz or the exact load, and why; linear buckling of the truss by hand at $30^\circ$;
counting negative pivots; which root of Crisfield's quadratic; limit point or bifurcation from
a computed tangent; why linear buckling fails for a snap-through.}

\framesub{Problems}
\draftnote{Planned: Ritz and finite elements for the column (convergence from above);
the truss as two bar elements, tangent stiffness assembled by hand, $\mat{K}_0$, $\mat{K}_\sigma$
and the linear buckling loads at $30^\circ$ and $65^\circ$; the extended system for a critical
point and its Jacobian.}

\framesub{Computer exercises}
\draftnote{Planned, all on the two-bar truss of the Chapter~2 session, with known answers
($EA = L = 1$): $\varphi = 30^\circ$, limit load $0.0481$, linear buckling $0.250$;
$\varphi = 65^\circ$, bifurcation $0.2434$, $\eta_2 = -0.524$, limit $0.2865$, linear buckling $0.324$.
(a)~Python: Riks--Crisfield from scratch on the two-element model, reusing the code of
Chapter~\ref{ch:st}. (b)~FEniCS: first steps, the truss as two interval elements in the plane,
energy in UFL, Newton, linear buckling with SLEPc, arc-length on top. (c)~Cast3M: the same truss,
\texttt{KSIG} + \texttt{VIBR}, then \texttt{PASAPAS} with arc-length; branch switching with a
sideways load.}

\begin{chapterbib}{99}
\bibitem{nc-keller1977} H.~B.~Keller, Numerical solution of bifurcation and nonlinear eigenvalue problems, in P.~H.~Rabinowitz (ed.), \emph{Applications of Bifurcation Theory}, Academic Press, New York, 1977, 359--384.
\bibitem{nc-riks1979} E.~Riks, An incremental approach to the solution of snapping and buckling problems, \emph{International Journal of Solids and Structures} 15 (1979) 529--551.
\bibitem{nc-crisfield1981} M.~A.~Crisfield, A fast incremental/iterative solution procedure that handles snap-through, \emph{Computers \& Structures} 13 (1981) 55--62.
\bibitem{nc-allgower2003} E.~L.~Allgower and K.~Georg, \emph{Introduction to Numerical Continuation Methods}, SIAM Classics in Applied Mathematics, Philadelphia, 2003.
\bibitem{nc-seydel2010} R.~Seydel, \emph{Practical Bifurcation and Stability Analysis}, 3rd ed., Springer, New York, 2010.
\bibitem{nc-cochelin2007} B.~Cochelin, N.~Damil and M.~Potier-Ferry, \emph{M\'ethode asymptotique num\'erique}, Herm\`es-Lavoisier, Paris, 2007.
\bibitem{nc-logg2012} A.~Logg, K.-A.~Mardal and G.~N.~Wells (eds.), \emph{Automated Solution of Differential Equations by the Finite Element Method}, Springer, Berlin, 2012.
\bibitem{nc-triantafyllidis2011} N.~Triantafyllidis, \emph{Stability of Solids: from Structures to Materials}, lecture notes MEC563, \'Ecole polytechnique, Palaiseau, 2011.
\bibitem{nc-constantinescu2007} A.~Constantinescu and A.~Korsunsky, \emph{Elasticity with Mathematica: An Introduction to Continuum Mechanics and Linear Elasticity}, Cambridge University Press, 2007.
\end{chapterbib}
